\chapter{The repository}
\label{ch:repo}

\section{What is in it}

The repository is one Python package plus the files needed to install it,
one converted robot asset, and documentation. After the clean-up made in this
revision the tree looks like this:

\begin{lstlisting}[style=plain,numbers=none]
Luna_HiFi/                          <- git root, pyproject.toml lives here
  pyproject.toml                    how pip installs the package + the isaaclab.tasks entry point
  .gitignore                        keeps __pycache__, *.egg-info and LaTeX build files out of git
  README_RL.md                      the original template notes (superseded, kept for the excavation map)
  RUNBOOK.md                        the validated setup and daily commands
  luna_hifi_tasks/                  THE PYTHON PACKAGE
    __init__.py                     imports .tricycle -> registers the task
    tricycle/                       the pipeline-validation task (this manual)
      __init__.py                   gym.register("Luna-Tricycle-Direct", ...)
      tricycle.py                   the car as an MJCF string (source of the USD asset)
      tricycle_env_cfg.py           scene, physics, spaces, termination knobs
      tricycle_env.py               the environment class: actions, obs, reward, dones, reset
      agents/
        __init__.py                 empty; makes "agents" importable
        rsl_rl_ppo_cfg.py           PPO hyper-parameters for rsl_rl
    excavation/                     the future rover task (template, not covered here)
  assets/
    config.yaml                     converter settings that produced the USD (generated)
    .asset_hash                     converter cache key (generated)
    tricycle/
      tricycle.usda                 the file the env cfg points at
      payloads/base.usda            geometry (visible shapes) of the car
      payloads/robot.usda           Isaac robot/link/joint schema tags
      payloads/Physics/physics.usda rigid bodies, joints, masses, Newton collision tags
      payloads/Physics/mujoco.usda  MuJoCo-specific overrides (joint damping)
      payloads/Physics/physx.usda   PhysX-specific overrides (unused with Newton)
  docs/tricycle_manual/             this document (LaTeX source + PDF)
\end{lstlisting}

Three kinds of file live here and it helps to keep them apart:

\begin{itemize}[leftmargin=1.6em]
  \item \textbf{Code you write.} Everything in \srcpath{luna_hifi_tasks/} and
        \code{pyproject.toml}.
  \item \textbf{Generated files you keep.} Everything in \srcpath{assets/}. They
        are produced by Isaac Lab's MJCF-to-USD converter from
        \code{tricycle.py} and are committed because the converter needs the
        Isaac Sim pip package, which not every machine has. When
        \code{tricycle.py} changes, these must be regenerated
        (Chapter~\ref{ch:runbook}).
  \item \textbf{Generated files you do not keep.} \srcpath{__pycache__/} folders
        and \srcpath{luna_hifi_tasks.egg-info/}. They are rebuilt by Python and
        pip every time and are now ignored by git.
\end{itemize}

\section{When does each file run?}
\label{sec:when}

A useful way to understand the package is to ask \emph{when} each file is
executed. There are four moments.

\begin{description}[leftmargin=1.6em,font=\bfseries]
  \item[Install time] \code{pip install -e .} reads \code{pyproject.toml}. It
    records two things in the Python environment: that the folder
    \code{luna_hifi_tasks} is importable, and that the string
    \code{"luna_hifi_tasks"} is registered under the entry-point group
    \code{isaaclab.tasks}. Nothing in the package is executed yet.
  \item[Start-up time] When you run \code{isaaclab train ...}, Isaac Lab's
    command-line tool asks Python for every entry point in the
    \code{isaaclab.tasks} group and imports each one. Importing
    \code{luna_hifi_tasks} runs \srcpath{luna_hifi_tasks/__init__.py}, which
    imports \srcpath{tricycle/__init__.py}, which calls \code{gym.register}.
    From now on Gymnasium knows the name \code{Luna-Tricycle-Direct} and
    which classes it maps to. \srcpath{tricycle/__init__.py} also imports
    \code{agents}, so \srcpath{agents/__init__.py} runs (it is empty).
  \item[Configuration time] Isaac Lab looks up the registered task, imports
    \code{tricycle_env_cfg.py} and \srcpath{agents/rsl_rl_ppo_cfg.py}, creates
    one \code{TricycleEnvCfg} and one \code{TricyclePPORunnerCfg}, and applies
    command-line overrides such as \code{--num_envs 4} to them. Importing
    \code{tricycle_env_cfg.py} imports \code{tricycle.py} for the constants
    it needs (\code{CHASSIS_Z}, joint names).
  \item[Run time] Isaac Lab imports \code{tricycle_env.py}, builds a
    \code{TricycleEnv} from the config, and hands it to \code{rsl_rl}. From
    then on the only code of yours that runs is the handful of methods in
    \code{TricycleEnv}, called thousands of times per second.
\end{description}

Chapter~\ref{ch:pipeline} follows this sequence in detail.

\section{Tricycle and excavation are independent}
\label{sec:independent}

The two task folders share nothing at import time. The evidence:

\begin{itemize}[leftmargin=1.6em]
  \item \srcpath{luna_hifi_tasks/__init__.py} contains a single line:
        \code{from . import tricycle}. The \code{excavation} package is never
        imported, so none of its code runs and none of its \code{VERIFY}
        markers can affect a tricycle run.
  \item No file in \srcpath{tricycle/} mentions \code{excavation}, and no file in
        \srcpath{excavation/} mentions \code{tricycle}. Each has its own
        \code{__init__.py}, its own env, its own cfg and its own
        \srcpath{agents/}.
  \item Until this revision, \srcpath{tricycle/soil.py} was a byte-for-byte copy
        of \srcpath{excavation/soil.py} that nothing in \srcpath{tricycle/}
        imported. The tricycle copy has been removed; the excavation copy is
        untouched (Appendix~\ref{app:changes}).
\end{itemize}

The only thing they share is the Isaac Lab machinery they both sit on. When
the excavation task is finished, adding \code{from . import excavation} to
\srcpath{luna_hifi_tasks/__init__.py} will register it next to the tricycle.

\section{The runbook and this manual}

\code{RUNBOOK.md} is the operational document: what to type, in which
window, and what the error messages mean. This manual is the explanatory
document: why the code is the way it is. Chapter~\ref{ch:runbook} goes through
the runbook's commands flag by flag so the two stay connected.
